\documentclass[a4paper]{article}

% macOS: use fontset=mac to avoid glyph gaps in 昇/騰 etc.
\usepackage[fontset=mac]{ctex}
\usepackage{amsmath, amssymb}
\usepackage{graphicx}
\usepackage[margin=2.5cm]{geometry}
\usepackage[most]{tcolorbox}
\usepackage{etoolbox}
\usepackage{listings}
\usepackage{hyperref}
\usepackage{booktabs}
\usepackage{subcaption}
\usepackage{float}
\usepackage{tikz}
\usetikzlibrary{positioning, arrows.meta, shapes.geometric, calc, chains}

% Highlight boxes
\newtcolorbox{knowledgebox}[1]{
    enhanced, colback=blue!5!white, colframe=blue!75!black, colbacktitle=blue!75!black,
    coltitle=white, fonttitle=\bfseries, title=#1, attach boxed title to top left={yshift=-2mm, xshift=2mm},
    boxrule=1pt, sharp corners
}

\newtcolorbox{importantbox}[1]{
    enhanced, colback=yellow!10!white, colframe=yellow!80!black, colbacktitle=yellow!80!black,
    coltitle=black, fonttitle=\bfseries, title=#1, sharp corners
}

\newtcolorbox{warningbox}[1]{
    enhanced, colback=red!5!white, colframe=red!75!black, colbacktitle=red!75!black,
    coltitle=white, fonttitle=\bfseries, title=#1, sharp corners
}

% Listings
\lstset{
    language=Python,
    basicstyle=\ttfamily\small,
    keywordstyle=\color{blue},
    stringstyle=\color{red!60!black},
    commentstyle=\color{green!60!black},
    breaklines=true,
    frame=single,
    numbers=left,
    numberstyle=\tiny\color{gray},
    captionpos=b,
    extendedchars=true,
    inputencoding=utf8
}

% Front-page metadata
\newcommand{\notetitle}{ai infra职业规划-2026年7月18日-学历和实习都顶级拉满的一位同学}
\newcommand{\noteauthors}{五道口纳什 \& Codex}
\newcommand{\notedate}{2026年7月19日}
\newcommand{\videochannel}{五道口纳什}
\newcommand{\videopublishdate}{2026年7月18日}
\newcommand{\videoduration}{25分15秒}
\newcommand{\videourl}{https://www.bilibili.com/video/BV18VKB6zE51/}
\newcommand{\videocoverpath}{cover.jpg}

\begin{document}

\begin{titlepage}
\centering
{\Large 课程笔记\par}
\vspace{1.2cm}
{\huge\bfseries \notetitle\par}
\vspace{0.8cm}
{\large \noteauthors\par}
\vspace{0.3cm}
{\large \notedate\par}
\vspace{1.2cm}

\ifx\videocoverpath\empty
{\small 请在生成笔记时填入视频封面图的本地路径。\par}
\else
\includegraphics[width=0.82\textwidth,height=0.45\textheight,keepaspectratio]{\videocoverpath}\par
\fi

\vfill
\begin{tcolorbox}[width=0.9\textwidth, colback=black!2!white, colframe=black!60, sharp corners]
\textbf{视频作者/频道}：\videochannel\par
\textbf{发布日期}：\videopublishdate\par
\textbf{视频时长}：\videoduration\par
\textbf{视频链接}：
\ifx\videourl\empty
未填写
\else
\href{\videourl}{\nolinkurl{\videourl}}
\fi
\end{tcolorbox}
\end{titlepage}

\tableofcontents
\newpage

% ============================================================
\section{咨询背景与学员画像}

本次咨询为 AI Infra 方向一对一职业规划。学员背景突出：来自国内顶级高校（985），在头部 GPU 厂商完成两段实习，学历与实习经历均属校招候选人中的前列。

学员之前在一家芯片公司做激活算子开发（基于 TensorFlow 框架），后经咨询老师建议跳槽到当前公司。学员对跳槽选择表示「特别十分满意」，简历也比之前充实不少。咨询围绕四个核心问题展开：

\begin{enumerate}
\item 简历中同一家公司的实习经历如何呈现：摘要式 vs 细目式？
\item 技术方向选择：继续深耕算子，还是借组内机会转向 RL / 端到端优化？
\item 是否需要回学校发论文，还是仅凭实习经历即可？
\item 求职目标选择：大厂 AI Infra 核心组 vs AI 独角兽？
\end{enumerate}

\begin{knowledgebox}{学员画像速览}
\begin{itemize}
\item \textbf{学历}：985 本硕，导师方向为存算一体（边缘端，偏学术）
\item \textbf{实习}：两段 GPU 厂商经历，覆盖算子开发、RL 后训练、推理框架调试
\item \textbf{技术栈}：CUDA 算子、CURL、PyTorch 框架、TP 通信、matmul 优化
\item \textbf{当前定位}：算子组正式员工，同时参与 RL 组和推理组的跨组项目
\item \textbf{求职窗口}：2027 届秋招（2026 年 8--9 月开始投递）
\end{itemize}
\end{knowledgebox}

\begin{figure}[H]
\centering
\includegraphics[width=0.75\textwidth]{figures/talk_180s.jpg}
\caption{咨询开场——老师查看学员简历 Word 文档，确认前期咨询的跳槽建议落实情况\protect\footnotemark}
\end{figure}
\footnotetext{视频画面时间区间：00:00:00--00:03:00。}

\subsection{本章小结}
学员基础扎实但面临方向选择：算子路线已达标，RL/端到端优化是 2026 年新增增长点。简历呈现、技术方向、论文策略、求职目标四个问题贯穿全篇。

% ============================================================
\section{简历结构：摘要式 vs 细目式}

学员在同一家公司有两段实习，简历上分别以两种风格呈现：

\begin{itemize}
\item \textbf{版本一（摘要式）}：三个高概括度条目——稳定性修复、Debug、RL 后训练，每条带明确工作背景和量化结果
\item \textbf{版本二（细目式）}：逐条罗列技术细节，包含 M\&K 配置、预编译目录、通信队列状态等底层描述
\end{itemize}

老师的判断非常直接：\textbf{版本二「太像一个实验记录，不太像项目经理」}。面试官阅读简历时首先需要的是\textbf{工作背景}——这个工作为什么要做，而不是直接跳到技术细节。版本一的写法让面试官能快速理解「你在解决什么问题」，版本二虽然技术指标详实，但缺乏前因后果的叙事线。

\begin{importantbox}{简历写作核心原则}
\begin{itemize}
\item \textbf{先给背景，再给动作}：每条经历必须回答「为什么做这件事」+「做了什么」+「结果如何」
\item \textbf{面试时再展开细节}：简历是敲门砖，不是技术报告。保留主干，面试中口头补充技术深度
\item \textbf{避免实验记录风}：逐条罗列配置调参就像实验 logbook，面试官无法判断你的工作价值和独立思考能力
\end{itemize}
\end{importantbox}

另外，老师建议在每条经历前加\textbf{粗体小标题}（十个字以内），让面试官第一眼就能抓住重点。经历排序上，将更充实的段落放在前面。

\subsection{本章小结}
简历采用摘要式+背景驱动的写法，保留主干、预留细节在面试中展开。加粗体小标题、调整段落顺序可显著提升可读性。

% ============================================================
\section{技术经历深度挖掘}

学员在咨询中展示了三个有代表性的技术经历，老师逐一评估表达方式和面试呈现策略。

\subsection{MatMul 预编译优化}

学员参与了一个 matmul 算子的优化工作，核心思路是\textbf{预编译目录}：将不同 M/K/CONFIG 的编译结果预先缓存到目录中，模型直接走时无需重复编译，节省 JIT 编译开销。

老师指出这个工作偏向\textbf{研究性/调试性任务}——不是从零开发，而是同事把设置搞错后学员去 debug 并优化。简历中只需按摘要式写法保留一条即可，不需要展开具体的 CONFIG 列表和参数调优细节。

\begin{warningbox}{调参 vs 开发——面试官视角的区分}
面试官能分辨「调参/修复已有代码」和「从零设计并实现一个功能」。前者的技术深度有限，简历中占比不宜过大。应侧重展示你在 debug 过程中\textbf{分析问题的思路}，而非最终调到什么参数值。
\end{warningbox}

\subsection{TP 通信死锁 Debug}

这是一个更具技术深度的经历。学员排查了一个 Tensor Parallel 通信中的 hang 问题：

\begin{itemize}
\item 两张 GPU 通过共享队列进行 TP 通信
\item 队列的状态改变（state change）是触发下一阶段通信的条件
\item 由于 FTC（Fused Token Communication）传入的数据偏小，导致状态误判——系统认为数据未到齐，陷入死循环等待
\item 最终整个进程超时被 KO
\end{itemize}

老师在听完表述后给出了积极评价：「我听起来还好，你对这个算子还是比较熟悉的。」这段经历展示了学员对 TP 通信机制的底层理解，在面试中有较大加分空间——关键是讲清楚\textbf{根因分析链条}而非仅描述现象。

\subsection{训练-推理专家选择分歧}

学员发现推理阶段和训练阶段的 MoE 专家选择存在差异。他做了定性分析：即使面对相同输入，推理和训练调用的算子不同，导致的微小偏差在 ResNorm 中逐层累积，最终选出不同的专家。

老师追问是否有定量分析——学员坦诚仅做了定性分析。这段经历的价值在于展示了学员对 MoE 架构（Expert Parallelism + All-to-All 通信）的理解，但面试中应主动说明「做了定性分析，定量分析是下一步」。

\begin{figure}[H]
\centering
\includegraphics[width=0.75\textwidth]{figures/talk_600s.jpg}
\caption{技术经历深度讨论——老师逐一评估学员的三段 Debug/优化经历，给出面试表达建议\protect\footnotemark}
\end{figure}
\footnotetext{视频画面时间区间：00:04:00--00:10:00。}

\subsection{本章小结}
三段经历代表了 AI Infra 实习生的典型工作谱系：开发型优化（matmul）、底层通信 Debug（TP hang）、架构级问题分析（专家选择分歧）。面试表达关键是「讲清楚为什么→怎么做→发现什么」，而非堆砌技术参数。

% ============================================================
\section{算子开发三场景：Prefill / Decode / Mixed}

学员在算子开发中同时支持了三个推理场景：

\begin{itemize}
\item \textbf{Prefill（预填充）}：首次处理 prompt，计算密集，对算子吞吐敏感
\item \textbf{Decode（自回归解码）}：逐 token 生成，内存带宽敏感，batch 效率关键
\item \textbf{Mixed（混合调度）}：prefill + decode 并发，需要调度策略协调两类算子的资源竞争
\end{itemize}

学员关注的问题是：是否需要在简历中说明三个场景的实现差异？老师的建议是\textbf{值得写}——因为三个场景的算子实现确有差异（纯 prefill / 纯 decode / 混合各自的最优策略不同），而且学员额外设计了一个\textbf{fallback 路径}（当优化路径不可用时回退到正常情况），这展示了工程上的鲁棒性思维。

\begin{knowledgebox}{SyncNicOp 与同步开销}
学员提到的 SyncNicOp（Synchronized NIC Operation）是一种算子执行模式，允许单个算子同时处理 prefill 和 decode 请求。但其代价是在 CPU 上消耗大量同步等待时间（等待通信完成），三个场景的优化正是要减少这种同步开销。学员的 fallback 机制在优化路径失效时回退到标准 SyncNicOp，保证功能正确性。
\end{knowledgebox}

\subsection{本章小结}
Prefill/Decode/Mixed 三场景的算子优化是区分初级和资深算子工程师的关键标志。加上 fallback 路径的设计，这段经历在简历中应作为独立条目写出。

\clearpage

% ============================================================
\section{职业方向：从算子到端到端优化}

这是本次咨询最核心的方向性讨论。学员的问题是：\textbf{继续深耕算子，还是利用组内机会转向 RL / 推理框架的端到端优化？}

老师的回答分两层：

\textbf{第一层：算子路线已经达标。}学员两段实习的主线都是算子开发，技术深度和产出已经达到秋招找到好工作的要求，不需要继续在算子这个单一维度上叠加。

\textbf{第二层：端到端优化是视野升级。}老师指出一个关键认知——\textbf{端到端延迟并不全是算子决定的}。推理系统的延迟由多因素构成：
\begin{itemize}
\item 算子计算效率（学员当前的主线）
\item \textbf{通信开销}：TP/PP/EP 的 All-Reduce、All-to-All 延迟
\item \textbf{调度策略}：prefill-decode 混合调度的排队延迟
\item \textbf{训练框架同步}：RL 后训练中训练框架和推理框架之间的同步开销
\end{itemize}

只做算子的人看不到通信和调度层面的瓶颈，而\textbf{RL 后训练}恰好要求同时理解算子、通信、调度、训练框架——这是一个自然的能力扩展路径。

\begin{figure}[H]
\centering
\begin{tikzpicture}[
    node distance=0.6cm and 1.5cm,
    box/.style={rectangle, draw, minimum width=3.8cm, minimum height=0.75cm, align=center, font=\small, rounded corners=2pt},
    arrow/.style={thick,->,>=stealth}
]
\node [box, fill=blue!8] (op) {算子开发\\（学员已有）};
\node [box, fill=green!8, right=of op] (rl) {RL 后训练\\（当前机会）};
\node [box, fill=yellow!8, right=of rl] (e2e) {端到端优化\\（目标状态）};

\draw [arrow] (op) -- (rl) node[midway, above, font=\footnotesize] {跨组参与};
\draw [arrow] (rl) -- (e2e) node[midway, above, font=\footnotesize] {视野升级};

\node [box, fill=red!5, below=of rl, yshift=-0.4cm] (io) {IO Infra\\（2026 新兴）};
\draw [arrow, dashed] (rl) -- (io) node[midway, right, font=\footnotesize] {自然延伸};

\node [box, fill=orange!8, below=of op, yshift=-0.4cm] (comm) {通信/调度优化\\（视野补全）};
\draw [arrow, dashed] (op) -- (comm) node[midway, left, font=\footnotesize] {短板补齐};
\end{tikzpicture}
\caption{算子工程师的能力扩展路径：从单点优化到端到端系统视野}
\end{figure}

\begin{importantbox}{能力扩展策略}
\textbf{不要放弃算子}——算子深度是竞争壁垒。在此基础之上，借 RL 后训练的跨组机会\textbf{扩展系统视野}，理解通信调度、训练-推理同步、端到端延迟分解。这样秋招时你的故事是「算子专家 + 系统视角」，而非「只会写 kernel 的人」。
\end{importantbox}

\subsection{本章小结}
算子路线已达标，应借 RL 后训练机会补齐通信调度和端到端瓶颈分析能力。从「算子工程师」升级为「AI Infra 系统工程师」是秋招市场的差异化定位。

% ============================================================
\section{IO Infra：2026 年新兴赛道}

老师在讨论学员 RL 后训练工作时，特别提到了 \textbf{IO Infra} 这一新兴方向：

\begin{itemize}
\item 2025 年之前，IO Infra 并非独立岗位——几乎没听说过有人专门做 IO Infra
\item 从 2026 年开始，无论是社招市场还是学员反馈，IO Infra 的招人需求\textbf{越来越明确}
\item RL 后训练中，checkpoint 读写、数据集加载、分布式训练的数据管线都是 IO 瓶颈
\end{itemize}

学员目前的 RL 后训练工作恰好涉及 IO Infra 相关内容（简历中已列出）。老师建议\textbf{不要把这个方向当副线处理}——IO Infra 正处在从无到有的需求爆发期，提前卡位有先发优势。

\begin{knowledgebox}{IO Infra 为什么在 2026 年爆发？}
RL 后训练的规模化（RLHF/DPO/GRPO）带来了前所未有的 I/O 压力：每个训练 step 需要同时读写多个模型的 checkpoint、加载偏好数据集、记录 reward 信号。传统 Infra 团队往往把 I/O 当次要问题，但当模型规模达到 100B+ 且训练集群超过千卡时，I/O 成为\textbf{首要瓶颈}。
\end{knowledgebox}

\subsection{本章小结}
IO Infra 是 2026 年从 0 到 1 的新赛道，需求明确且供给稀缺。学员当前的 RL 后训练经历恰好覆盖这一方向，值得在简历和面试中重点突出。

\clearpage

% ============================================================
\section{Blackwell 架构迁移与 TMA}

学员提出：正职同事有 NVIDIA Blackwell（B 卡）的测试权限，是否可以合作把 Hopper（H 卡）上的算子迁移到 B 卡上？这段经历会不会是一个加分项？

老师的回答：\textbf{可以，但不要局限于参数调优}。Blackwell 架构相比 Hopper 的核心变化之一是 \textbf{TMA（Tensor Memory Access）}——一种新的硬件加速数据传输机制，可以显著降低 tensor 在 HBM 和 SMEM 之间的搬运延迟。迁移算子到 B 卡时，如果仅做编译参数调整（M/K 改一改），面试价值有限；但如果深入利用 TMA 做数据搬运优化（比如用 TMA 替代传统的 `cp.async`），则展示了对新硬件架构的深度理解。

\begin{warningbox}{Blackwell 迁移的常见误区}
\begin{itemize}
\item \textbf{误区}：以为迁移就是改 M/K/CONFIG 让 kernel 能在 B 卡上跑起来
\item \textbf{正确姿势}：理解 B 卡架构新特性（TMA、SM 数量变化、HBM 带宽提升），\textbf{针对新特性重新设计数据流}，而非仅做兼容性适配
\item \textbf{面试价值排序}：利用 TMA 设计新数据流水线 $>$ 调参适配 $>$ 仅跑通
\end{itemize}
\end{warningbox}

\subsection{本章小结}
B 卡迁移的经历有加分空间，但加分大小取决于工作深度。利用 TMA 做数据搬运重设计是正确打开方式，仅做参数适配面试价值不大。

% ============================================================
\section{技术栈准备策略}

学员提出了两个具体的技术准备问题：

\paragraph{CUDA/CURL 算子熟练度}学员表示「现在也会写，但肯定不是写的特别熟练」。老师建议\textbf{不必追求完美娴熟}，而是通过 RL 后训练和 B 卡迁移中的实际需求来\textbf{以战代练}——在具体优化实践中磨练 CURL 能力，比单纯刷教程更高效。

\paragraph{LeetCode 算法}学员之前刷过力扣，但近期松懈。老师建议 8 月份开始集中恢复，因为 AI Infra 公司的面试中算法题仍是必考项。但要注意\textbf{精力分配的优先级}：RL 后训练 + IO Infra 工作 $>$ B 卡迁移 $>$ LeetCode 刷题。面试前一个月集中刷即可，平常用 10\%--20\% 精力维持手感。

另外，老师特别建议在简历中\textbf{加粗 vLLM 关键词}，因为学员对 vLLM 框架比较熟悉，而在当前国内市场 vLLM 占明显份额优势（见第 10 节）。

\subsection{本章小结}
CURL 以战代练、LeetCode 8 月集中恢复。vLLM 关键词在简历中加粗突出。三条线的优先级：实际工作深度 $>$ 广度标签 $>$ 刷题准备。


% ============================================================
\section{求职策略与公司选择}

学员的核心问题是：\textbf{投大厂 AI Infra 核心组，还是投 AI 独角兽（DeepSeek、智谱、MiniMax 等）？}

老师的策略是一个词：\textbf{两边都投}。

\begin{importantbox}{求职投递策略}
\begin{enumerate}
\item \textbf{独角兽的门槛不低反高}：DeepSeek、智谱等公司对 AI Infra 的要求可能比大厂更高，因为它们对 Infra 的依赖更深
\item \textbf{大厂核心组 vs 边缘组需分辨}：有些大厂的 AI Infra 组仅是去支持业务，并非核心研发。优先瞄准做底层框架/训练系统的核心组
\item \textbf{用 offer 互相 leverage}：两边的 offer 可以用来在另一边谈薪资，这是校招中常见的议价手段
\end{enumerate}
\end{importantbox}

学员目前目标明确，8 月集中准备、9 月开始投递。老师特别提醒：

\begin{itemize}
\item 简历中把\textbf{某国产 GPU 公司}的名称亮出来——不需要遮遮掩掩，头部 GPU 厂商的实习经历在市场上是硬通货
\item 经历段落前加\textbf{粗体总标题}，并调整两段经历的先后顺序，让更充实的经历靠前
\item 简历制式化后先投一批中等目标「练手」，根据反馈调整，\textbf{不要海投}——有针对性地投 15--20 家，每家定制简历侧重点
\end{itemize}

\begin{figure}[H]
\centering
\includegraphics[width=0.75\textwidth]{figures/talk_1200s.jpg}
\caption{求职策略讨论——老师建议两边都投，用 offer 互相 leverage\protect\footnotemark}
\end{figure}
\footnotetext{视频画面时间区间：00:20:00--00:22:00。}

\subsection{本章小结}
大厂核心组和 AI 独角兽两边都投，用 offer 议价。简历亮出公司品牌、定制化投递 15--20 家，拒绝海投。

% ============================================================
\section{vLLM vs SGLang 市场现状}

学员在咨询结尾问了一个务实的问题：\textbf{vLLM 和 SGLang 应该重点选哪个？}

老师给出了一个具体数字：\textbf{自己的统计大概是七比三}——vLLM 占 7 成，SGLang 占 3 成。

\begin{importantbox}{框架选择建议}
\begin{itemize}
\item \textbf{选多数派（vLLM）}：给自己更多的选择空间。市场占有率高意味着更多公司在使用、更多岗位需要。
\item \textbf{简历中突出 vLLM}：在项目经历中加粗「vLLM」关键词，如果确实对 vLLM 源码和架构比较熟悉，可以适当在经历描述中提一提。
\item \textbf{不放弃 SGLang}：SGLang 在特定场景（如结构化生成）有独特优势，了解即可，不用深入。
\end{itemize}
\end{importantbox}

这个七三开的判断基于老师自己的行业观察（学员 + 社招市场），并非严格统计，但方向性明确——vLLM 目前是国内 LLM 推理框架的事实标准。

\subsection{本章小结}
vLLM 市场份额约 7 成，SGLang 约 3 成。求职策略选多数派，简历中加粗 vLLM 关键词。同时了解 SGLang 的结构化生成优势作为差异化知识储备。

\clearpage

% ============================================================
\section{总结与延伸}

\subsection{讲师结尾讨论}

老师在咨询结尾做了简要总结，核心要点：

\begin{itemize}
\item 学员目标明确、想法清晰，基础已经很好，只需在细节上打磨
\item 简历格式优化（粗体标题、经历排序、公司品牌亮出）是低投入高回报的最后一步
\item 8 月集中精力在 RL 后训练 + B 卡迁移上，9 月开始投递
\item 尾声以「期待你的好消息」收尾，鼓励学员高效执行
\end{itemize}

\subsection{结构化蒸馏}

本次咨询虽然是一对一对话，但其中蕴含的 AI Infra 职业规划逻辑具有普适性。以下是核心主张的提炼：

\begin{enumerate}

\item \textbf{简历的本质是「叙事」而非「清单」}：面试官阅读简历的时间通常不超过 30 秒。能在 30 秒内让面试官理解「你做了什么 + 为什么做」的简历，远胜于堆满技术参数的实验记录。摘要式写法 + 粗体关键词 + 背景驱动叙事是 AI Infra 简历的最佳实践。

\item \textbf{算子深度是护城河，系统视野是天花板}：纯算子工程师的市场需求存在，但增长空间有限。从算子出发，借 RL 后训练或推理框架的跨组机会补齐通信调度、端到端瓶颈分析能力，是从「执行者」到「设计者」的关键一跃。

\item \textbf{2026 年 AI Infra 的两个结构性机会}：
\begin{itemize}
\item \textbf{RL 后训练}：post-training 的规模化落地带来了大量 Infra 需求（训练-推理同步、reward pipeline、多模型 checkpoint 管理）
\item \textbf{IO Infra}：从无到有的新岗位，供给稀缺，需求在 2026 年明确爆发
\end{itemize}

\item \textbf{论文不是必需品}：在 AI Infra 工程岗位，顶级公司的实习经历比学术论文更有说服力。除非导师方向直接相关且不影响实习，否则不必「回学校发论文」。存算一体等偏学术的方向，对工程岗位的加分有限。

\item \textbf{新硬件 = 新机遇}：Blackwell 架构的普及意味着大量 Hopper 时代的算子需要重新设计。率先掌握 TMA 等新特性的工程师将在 2026--2027 年秋招中占据结构性优势。

\item \textbf{选框架就是选市场}：vLLM 以 7:3 的优势领先 SGLang，求职阶段选多数派是风险最低的策略。入职后根据实际需求再深入其他框架。

\end{enumerate}

\subsection{开放问题与下一步}

\begin{itemize}
\item RL 后训练中的 IO Infra 具体包含哪些子方向（data pipeline、checkpoint 管理、分布式文件系统优化）？哪些与学员现有工作最匹配？
\item Blackwell TMA 相比 Hopper 的 `cp.async` 在编程模型上有哪些具体差异？迁移一个典型 matmul kernel 需要多大改动？
\item vLLM 的 PagedAttention 和 SGLang 的 RadixAttention 在设计哲学上有何本质区别？各适合什么场景？
\item 2026 年秋招市场对 AI Infra 岗位的薪资水平和 HC 数量有何趋势？大厂和独角兽的 offer 结构有何差异？
\end{itemize}

\subsection{拓展阅读}

\begin{itemize}
\item NVIDIA Blackwell Architecture Technical Brief —— 官方白皮书，了解 TMA、FP4/FP6、第二代 Transformer Engine 等架构新特性
\item vLLM 官方文档：\href{https://docs.vllm.ai}{https://docs.vllm.ai} —— PagedAttention 设计原理和最新进展
\item SGLang 官方文档：\href{https://sglang.readthedocs.io}{https://sglang.readthedocs.io} —— RadixAttention 与结构化生成
\end{itemize}

\end{document}
